\uuid{Jiv8}
\chapitre{Statistique}
\niveau{L2}
\module{Analyse de données}
\sousChapitre{Tests d'hypothèses, intervalle de confiance}
\titre{Lois pour les statistiques}
\theme{loi normale, loi du chi2, loi de Student, loi de Fisher}

\auteur{}
\datecreate{2022-08-25}
\organisation{AMSCC}
\difficulte{}
\contenu{
\texte{$\mathcal N(\mu,\sigma)$ désigne la loi normale de moyenne $\mu$ et d’écart-type $\sigma$. Les quantiles sont à gauche : $q_{\nu;\gamma}$ pour la loi $\chi^2(\nu)$ ; $t_{\nu;\gamma}$ pour la loi $St(\nu)$ ; $F_{\nu_1,\nu_2;\gamma}$ pour la loi $F(\nu_1,\nu_2)$, avec une probabilité $\gamma$ à gauche du quantile.}
\texte{Les variables $U_1,\ldots,U_9$ sont indépendantes, de loi $\mathcal N(0,1)$.}
\begin{enumerate}
\item \question{Déterminer la loi de $X=\sum_{i=1}^9U_i^2$ et $x$ tel que $\mathbb P(X>x)=0{,}05$.}
\reponse{On a $X\sim\chi^2(9)$ et $x\simeq16{,}92$.}
\item \question{Soit $Y\sim\mathcal N(10,3)$ indépendant de $X$. Déterminer la loi de $Z=\frac{Y-10}{\sqrt X}$ et $z$ tel que $\mathbb P(Z>z)=0{,}05$.}
\reponse{On a $Z=\frac{\frac{Y-10}{3}}{\sqrt{\frac{X}{9}}}$ donc par définition, $Z$ suit une loi de Student $St(9)$ et $z\simeq1{,}833$.}
\item \question{Soit $V\sim\chi^2(3)$ indépendant de $X$. Déterminer les lois de $W_1=\frac{X}{3V}$ et $W_2=\frac{3V}{X}$, puis $w_2$ tel que $\mathbb P(W_2>w_2)=0{,}05$.}
\reponse{On a
\[
W_1=\frac{\frac X9}{\frac V3},
\]
donc $W_1$ suit une loi de Fisher $F(9,3)$. De même,
\[
W_2=\frac{\frac V3}{\frac X9},
\]
donc $W_2$ suit une loi de Fisher $F(3,9)$. On a
\[
w_2=F_{3,9;0{,}95}\simeq3{,}863.
\]}
\end{enumerate}
}
